\section{Related work}
\label{sec:related}

Two literatures bear on this question, both mature and well instrumented. Their
intersection is empty.

\paragraph{Growing and regenerating neural cellular automata.}
\citet{mordvintsev2020growing} showed that a shared local update rule can grow a
target morphology from a single seed and regenerate it after tissue loss, and
supplied the training recipe --- sample pool, stochastic updates, alive masking
--- that this work builds on directly. \citet{randazzo2020selfclassifying}
extended it to a per-cell classification task in which cells reach agreement on a
global property of the pattern they collectively form, and established two
choices we adopt: per-cell supervision rather than supervision of an aggregate,
and $L_2$ against a one-hot target rather than cross-entropy, since the latter
drives per-cell logits upward without bound and prevents the automaton settling.
Neither work involves more than one organism, and neither involves a learned
communication channel.

\citet{variengien2021towards} is the closest feasibility precedent for the
present question: an NCA controlling a cart-pole exhibits a developmental phase,
regeneration after damage, and robustness to input deletion. That establishes
that NCA-borne \emph{computation} can survive structural damage. It is
single-agent, and there is no protocol to lose.

\paragraph{Multiple organisms on a shared substrate.}
\citet{stovold2025identity} adds an identity channel that lets NCA organisms
remain stable in close proximity, framed explicitly as groundwork for studying
cellular-level social interaction --- but there is no signalling, no task, and no
protocol. Sakana AI's Petri Dish NCA \citep{sakana2026pdnca} lets many
independently parameterised NCA agents coexist, compete, and continually adapt in
one environment, and \citet{berdica2026evolving} scale that substrate with
population-based training toward open-ended discovery. Interaction there runs
through continuous attack and defence channels; there is no discrete vocabulary
and nothing that could be called a convention.

\paragraph{The nearest miss, and why it is still a miss.}
\citet{singh2026communication} is the closest existing work and deserves to be
separated from the others carefully, because a superficial reading makes it look
like this paper has already been written. That work studies NCA cells reaching
\emph{majority consensus} on an unobservable global density, and introduces
``languages'' as subpopulations separated by tunable translation barriers,
finding that protocol heterogeneity slows consensus and that rules trained under
varied protocols are more robust than homogeneously trained ones.

The consensus machinery genuinely overlaps with ours: our readout
(Section~\ref{sec:readout}) is a per-cell vote pooled into a collective decision,
and that is the same family of mechanism. We do not claim otherwise. The
divergence is on two other axes. First, the ``languages'' there are
\emph{imposed} --- a translation barrier is a designed parameter, not a
convention the agents invent --- whereas the object of study here is a protocol
that emerges from a referential game with no prescribed mapping from referents to
symbols. Second, and more decisively, nothing in that setting is damaged. The
question of whether a protocol survives the destruction and regrowth of the
substrate computing it does not arise.

\paragraph{Emergent communication.}
The Lewis signalling game \citep{lewis1969convention} is the standard setting,
and \citet{lowe2019pitfalls} is the standard we hold ourselves to: agents can
appear to communicate while doing nothing of the kind, so positive signalling and
positive listening must be demonstrated causally rather than inferred from task
success. We adopt their interventions as mandatory and enumerate counterfactual
symbols exactly rather than sampling them.
\citet{chaabouni2020compositionality} found no correlation between the
compositionality of an emergent language and its ability to generalize, which is
why we report compositionality-adjacent measures beside accuracy and never as a
proxy for it. What none of this literature provides is a notion of \emph{damage}:
its agents are fixed-architecture networks, and there is no operation on them
corresponding to removing 30\% of the machine and letting it grow back.

One apparent match is a false positive worth naming, since a keyword search
surfaces it immediately. ESCELL \citep{chowdhury2020escell} runs a genuine Lewis
referential game and has ``cellular'' in its title, but the word refers to
biological cells in microscopy images; the agents are standard convolutional
networks.

\paragraph{The gap.}
No prior work combines (i) NCA-grown, damage-regenerating agents with (ii) a
learned, discrete, emergent protocol, tested for (iii) semantic survival under
substrate damage. Each nearest miss drops exactly one leg: \citet{singh2026communication}
has consensus and multiple protocols but imposes them and never damages anything;
\citet{stovold2025identity} has coexistence without signalling;
\citet{sakana2026pdnca} has interaction without a discrete vocabulary;
\citet{variengien2021towards} has damage-surviving computation without
communication; \citet{chowdhury2020escell} has the game without the substrate.

There is a second gap, which is about measurement rather than systems. The NCA
literature measures recovery as distance to a target morphology or as downstream
task score. The emergent-communication literature has never needed a definition
of recovery, because its agents cannot be structurally damaged. \emph{Semantic}
recovery --- whether the meaning came back, as distinct from whether the body or
the channel capacity did --- has no existing definition. Section~\ref{sec:protocol}
supplies one.
